\input{preamble}

\begin{document}

\title{Representation Theory}
\maketitle

\phantomsection
\label{section-phantom}

\tableofcontents

\section{Basic definitions}
\label{section-basic-definitions}

\begin{definition}
\label{definition-representation-lie-algebra}
A {\it representation of Lie algebra}
$\mathfrak{g}$ is a vector space $V$
and homomorphism
$$
\rho:\mathfrak{g} \to \text{End}(V)
$$
of Lie algebras, i.e.,
$$
\rho([x,y])=\rho(x)\rho(y)-\rho(y)\rho(x).
$$
\end{definition}

\begin{remark}
\label{remark-representations-are-modules}
Another name for a representation of a Lie
algebra $\mathfrak{g}$ is a
{\it $\mathfrak{g}$-module}. This is like the
abstract definition of an algebra
using a morphism. It's as if the morphism
lets us define a sort of product of
elements in the codomain by elements in the
domain.
\end{remark}

It is not possible to classify all
representations of a Lie algebras
$\mathfrak{g}$. But there is a theorem by
Weyl that says that if $\mathfrak{g}$
is finite-dimensional and (semi)simple, then
every finite-dimensional
representation of $\mathfrak{g}$
(i.e. a representation where $V$ is
finite-dimensional)
is isomorphic to a direct sum of irreducible
representations

\section{Representations of groups}
\label{section-representations-of-groups}

\noindent
Notes from
\href{http://verbit.ru/IMPA/CA-2026/slides-%
CA-2026-05.pdf}{
Misha's lecture}.

\begin{definition}
\label{definition-representation-group}
A {\it representation of a group}
is a vector space $V$ and a morphism
$\rho:G \to \text{GL}(V)$.
A {\it morphism} of $G$-representations
is a map of the underlying vector spaces
commuting
with the $G$-action, i.e.
$\varphi(gx)=g\varphi(x)$.
\end{definition}

\begin{definition}
\label{definition-irreducible-representation}
A group representation $G \to \text{GL}(V)$
is {\it irreducible} if $V$ has no
$G$-invariant subspaces,
i.e. if there are no subrepresentations.
It is called {\it semisimple} if it is a
direct sum
of irreducible representations.
\end{definition}

\begin{remark}
\label{remark-semis-iff-any-short-exact}
All representations are semismple if and only
if
any short exact sequence of
$G$-representations splits.
\end{remark}

\begin{proposition}
\label{proposition-finit-group-have-semis}
Let $\text{Rep}_k(G)$ be the category of
representations
of a finite group $G$ over $k$,
with $\text{char}(k)$ coprime with $|G|$.
Then any short exact sequence of
$G$-representations splits.
\end{proposition}

\begin{proof}
The idea is to construct a section of the
right-hand map
on a short exact sequence by averaging over
the group.
We need $\text{char}(k)$ coprime with $|G|$
so that
we can divide by $|G|$.
\end{proof}

This means that any finite-dimensional
representation
of finite group with $\text{char}(k)$ coprime
with $|G|$
is semisimple.

\begin{exercise}
\label{exercise-nonsemisimple-borel}
Consider the group $G$ of non-degenerate
upper triangular matrices $n \times n$.
Construct a non-semisimple representation of
$G$.
\end{exercise}

\begin{proof}
Assume $n\geq 2$ and let $G$ act on
$V=k^n$ by the standard representation.
Let $e_1,\ldots,e_n$ be the standard basis.
The line
$$
U=\langle e_1\rangle
$$
is $G$-invariant, since every upper
triangular matrix sends $e_1$ to a scalar
multiple of $e_1$.

We claim that $U$ has no $G$-invariant
complement. Suppose that such a complement
$W$ exists. Since $V=U\oplus W$, there is a
unique vector $w\in W$ whose image in
$V/U$ is the image of $e_2$. Thus
$$
w=e_2+ce_1
$$
for some $c\in k$.

Now consider the unipotent upper triangular
matrix
$$
u_t=
\begin{pmatrix}
1&t&0&\cdots&0\\
0&1&0&\cdots&0\\
0&0&1&\cdots&0\\
\vdots&\vdots&\vdots&\ddots&\vdots\\
0&0&0&\cdots&1
\end{pmatrix}.
$$
Since $W$ is assumed invariant, $u_t w$ must
belong to $W$. But
$$
u_t w=e_2+(c+t)e_1.
$$
This vector has the same image as $e_2$ in
$V/U$. By uniqueness of $w$, we must have
$u_t w=w$, hence $c+t=c$. Taking
$t\neq 0$ gives a contradiction.

Therefore the invariant subspace $U$ has no
$G$-invariant complement. Hence the standard
representation of $G$ on $k^n$ is not
semisimple.
\end{proof}

\begin{exercise}
\label{exercise-finite-field-nonsemisimple}
Construct a finite group $G$ and a
non-semisimple representation
$$
G \longrightarrow \text{GL}(V),
$$
where $V$ is a finite-dimensional vector
space over a finite field.
\end{exercise}

\begin{proof}
We may use the above example for any
nontrivial finite field - the group
of upper triangular matrices is finite
and we get a non-semisimple representation
in the exact same way.
\end{proof}

\begin{definition}
\label{definition-real-complex-quaternionic}
Let $V$ be a finite-dimensional irreducible
group representation over $\mathbb{R}$.
It is called {\it real} if
$\text{End}_G(V)=\mathbb{R}$, {\it complex}
if $\text{End}_G(V)=\mathbb{C}$, and
{\it quaternionic} if $\text{End}_G(V)$ is
the algebra of quaternions.
\end{definition}

\begin{remark}
\label{remark-real-complex-quaternionic}
This definition is not about the image of
the homomorphism
$G \to \text{GL}(V)$. It is about the
endomorphisms of $V$ which commute with the
whole $G$-action:
$$
\text{End}_G(V)=
\{T\in \text{End}_{\mathbb{R}}(V):
T\rho(g)=\rho(g)T\text{ for all }g\in G\}.
$$
Thus $\text{End}_G(V)$ is an algebra of
operators on $V$. For an irreducible real
representation, Schur's lemma says that this
is a division algebra over $\mathbb{R}$.
By Frobenius' theorem, it is isomorphic to
one of
$$
\mathbb{R},\qquad \mathbb{C},\qquad
\mathbb{H}.
$$
These are the real, complex, and
quaternionic cases.

In other words, $\text{End}_G(V)$ asks:
after $G$ acts, what linear operations on
$V$ are still invisible to $G$, because they
commute with everything $G$ does? If these
invisible operations form $\mathbb{H}$, the
representation is quaternionic.
\end{remark}

\begin{exercise}
\label{exercise-quaternionic-finite-group}
Find an example of a quaternionic
representation of a finite group.
\end{exercise}

\begin{proof}
Let $G=Q_8=\{\pm 1,\pm i,\pm j,\pm k\}$ be
the quaternion group, and let $V=\mathbb{H}$
be regarded as a real vector space. Let
$G$ act by left multiplication:
$$
\rho(q)(x)=qx.
$$
This is a real representation of $G$.

We first check that it is irreducible. Let
$0\neq W\subset V$ be a $G$-invariant real
subspace. Since $W$ is stable under left
multiplication by $i,j,k$, it is stable
under left multiplication by every
quaternion.
Choose $0\neq x\in W$. Since $x$ is
invertible in $\mathbb{H}$, left
multiplication by $x^{-1}$ shows that
$1\in W$. Hence $W=\mathbb{H}$.

Now compute the commuting endomorphisms.
For every $a\in \mathbb{H}$, right
multiplication
$$
R_a(x)=xa
$$
commutes with left multiplication by $G$.
Conversely, let $T\in \text{End}_G(V)$ and
put $a=T(1)$. Since $T$ commutes with left
multiplication by $i,j,k$, it commutes with
left multiplication by every quaternion.
Thus, for all $x\in\mathbb{H}$,
$$
T(x)=T(x\cdot 1)=xT(1)=xa.
$$
So every $G$-equivariant endomorphism is
some $R_a$. Therefore
$$
\text{End}_G(V)\cong \mathbb{H}^{op}
\cong \mathbb{H}.
$$
Thus $V$ is quaternionic.
\end{proof}

\begin{exercise}
\label{exercise-quaternionic-cyclic-17}
Find a quaternionic representation of a
cyclic group of order $17$, or prove that
such a representation cannot exist.
\end{exercise}

\begin{exercise}
\label{exercise-quaternionic-tetrahedron}
Find a quaternionic representation of the
group $G$ of isometries of a regular
tetrahedron, or prove that such a
representation cannot exist.
\end{exercise}

\begin{proof}
The full group of isometries of a regular
tetrahedron is isomorphic to $S_4$. We show
that it has no quaternionic irreducible real
representation.

We use the Frobenius-Schur indicator. If
$\chi$ is an irreducible complex character
of a finite group, then
$$
\nu(\chi)=
\frac{1}{|G|}\sum_{g\in G}\chi(g^2)
$$
is $1,0,-1$ according as the corresponding
real representation is of real, complex, or
quaternionic type.

Let the conjugacy classes of $S_4$ be
$$
C_1,\ C_2,\ C_{22},\ C_3,\ C_4
$$
with cycle types
$$
1,\quad (12),\quad (12)(34),\quad
(123),\quad (1234).
$$
Their sizes are $1,6,3,8,6$. Squaring sends
$C_1,C_2,C_{22}$ to $C_1$, sends $C_3$ to
$C_3$, and sends $C_4$ to $C_{22}$. Hence
$$
\nu(\chi)=
\frac{10\chi(C_1)+8\chi(C_3)
+6\chi(C_{22})}{24}.
$$

The irreducible characters of $S_4$ are
$$
\begin{array}{c|ccccc}
&C_1&C_2&C_{22}&C_3&C_4\\
\hline
\chi_1&1&1&1&1&1\\
\chi_{\mathrm{sgn}}&1&-1&1&1&-1\\
\chi_2&2&0&2&-1&0\\
\chi_3&3&1&-1&0&-1\\
\chi_3'&3&-1&-1&0&1
\end{array}
$$
Substituting these rows into the formula for
$\nu(\chi)$ gives $\nu(\chi)=1$ in every
case. Thus every irreducible real
representation of $S_4$ is of real type.
In particular, none is quaternionic.
\end{proof}

\begin{exercise}
\label{exercise-cyclic-rational-2d}
Let $V=\mathbb{Q}^2$ be an irreducible
$2$-dimensional representation of the cyclic
group $G=\mathbb{Z}/n$ over $\mathbb{Q}$.
\begin{enumerate}
\item Prove that the algebra
$\text{End}_G(V)$ of endomorphisms of this
representation is commutative and
$2$-dimensional.
\item Prove that $\text{End}_G(V)$ is a
number field isomorphic to
$\mathbb{Q}[\sqrt{a}]$ for some
$a \in \mathbb{Z}$.
\end{enumerate}
\end{exercise}

\begin{proof}
Let $g$ be a generator of $G$, and write
$$
A=\rho(g)\in \text{GL}_2(\mathbb{Q}).
$$
Geometrically, one may picture this as a
finite order rotation. Then an endomorphism
commuting with the representation should be
something which has the same symmetry: a
homothety, or the rotation itself, or a
linear combination of these two operations.
Over $\mathbb{Q}$ we should say this
algebraically rather than in terms of angles.

Since $G$ is cyclic, an endomorphism
$T\in\text{End}_{\mathbb{Q}}(V)$ is
$G$-equivariant if and only if it commutes
with $A$:
$$
\text{End}_G(V)=
\{T\in\text{End}_{\mathbb{Q}}(V):TA=AT\}.
$$
Thus the problem is to compute the
centralizer of the single matrix $A$.

The scalar matrices give the homotheties,
and $A$ itself is the algebraic replacement
for the rotation. Since the representation is
irreducible, $A$ cannot be scalar; otherwise
every line in $V$ would be $G$-invariant.
Therefore $I$ and $A$ are linearly
independent over $\mathbb{Q}$.

By Cayley-Hamilton, every polynomial in $A$
is of the form
$$
aI+bA,\qquad a,b\in\mathbb{Q}.
$$
So the algebra $\mathbb{Q}[A]$ has dimension
$2$ and is commutative. The remaining point
is to show that every matrix commuting with
$A$ already lies in $\mathbb{Q}[A]$.
\end{proof}





\section{Representations of SL(2,C)}
\label{section-sl2c}

\noindent
Klein classified nontrivial finite subgroups
of $\text{SL}_2(\mathbb{C})$
up to conjugacy:
\begin{itemize}
\item Cyclic groups $\mathcal{C}_m$, $m \geq
2$. $|\mathcal{C}_m|=m$,
$$
\mathcal{C}_m=\left\{\begin{pmatrix}
\zeta&0\\
0&\zeta^{-1}\end{pmatrix}:\zeta^m=1
\right\}.
$$
\item Dihedral groups $\mathcal{D}_m$, $m\geq
2$,
$$
\mathcal{D}_m=\left<\mathcal{C}_{2m},
\begin{pmatrix}
0&-1\\
1&0
\end{pmatrix}\right>.
$$
For example, $\mathcal{D}_2$ is the {\it
quaternion group}
$\left<\underbrace{\begin{pmatrix}
i&0\\
0&-i
\end{pmatrix}}_{i},
\underbrace{\begin{pmatrix}
0&-1\\
1&0
\end{pmatrix}}_{j}\right>$.

\item Binary tetrahedral group $\mathcal{T}$,
with $|\mathcal{T}|=24$.
\item Binary octahedral group, $\mathcal{O}$,
with $|\mathcal{O}|=48$.
\item Binary icosahedral group $\mathcal{I}$
with $|I|=120$.

These three are related to the rotation
groups of the
Platonic solids by the 2-to-1 covering
$\text{SU}_2 \to \text{SO}_3(\mathbb{R})$.
To classify the finite subgroups of
$\text{SL}_2(\mathbb{C})$:
first show that any such subgroup is
conjugate to a subgroup
of $\text{SU}_2$ so that its image is a
finite subgroup of
$\text{SO}_3(\mathbb{R})$. We can classify
those by considering their
action on the sphere and showing that they
have to be the
rotation group of one of the Platonic solids
or a dihedral group
or a cyclic group (related to planar
polygons).
\end{itemize}

\subsection{Unit quaternions and rotations}
\label{subsection-unit-quaternions-rotations}

\noindent
Notes from Conway's book on quaternions.

\begin{lemma}
\label{lemma-so3-simple-rotations}
$\text{SO}(3)$ is generated by simple
rotations, i.e. rotations which fix a
vector.
\end{lemma}

\begin{lemma}
\label{lemma-unit-quaternions-cover-so3}
Every element of $\text{SO}(3)$ has the form
$x \mapsto q^{-1}xq$ for some unit
quaternion $q$. The map
$$
q \longmapsto [q],
\qquad
[q](x)=q^{-1}xq
$$
is a $2$-to-$1$ homomorphism from the group
of unit quaternions to $\text{SO}(3)$.
\end{lemma}

\noindent
The imaginary quaternions
$$
\operatorname{Im}\mathbb{H}
=
\mathbb{R}i\oplus\mathbb{R}j
\oplus\mathbb{R}k
$$
are naturally $\mathbb{R}^3$. The unit
quaternions act on them by conjugation:
$$
\rho(\alpha)(\beta)
=
\alpha\beta\overline{\alpha}.
$$
Thus we get a homomorphism
$$
\rho:S^3 \longrightarrow
\text{SO}_3(\mathbb{R})
$$
with kernel $\{\pm 1\}$. This identifies the
unit quaternions with $\text{Spin}(3)$, and
also with $\text{SU}(2)$. The covering is
$2$-to-$1$ because opposite unit quaternions
give the same rotation.

\begin{definition}
\label{definition-binary-icosahedral}
The {\it binary icosahedral group} is the
preimage of the icosahedral rotation group
$I\cong A_5$ under the covering
$$
\text{SU}(2) \longrightarrow
\text{SO}_3(\mathbb{R}).
$$
\end{definition}

\begin{remark}
\label{remark-quaternions-binary-groups}
This is the bridge between the rotation
groups of the Platonic solids and the finite
subgroups of $\text{SU}(2)$, or equivalently
of $\text{SL}_2(\mathbb{C})$. The word
``binary'' means ``take the inverse image
under the double cover''. For example, the
quaternion group $Q_8$ is the binary
dihedral group $\mathcal{D}_2$.
\end{remark}

\noindent
At some point of the XX century it was
realised that this
classification is an instance of an ADE
classification.

Now we explain the (historically second
explanation of this) McKay
correspondence.
If $\Gamma$ is a nontrivial finite subgroup
of $\text{SL}_2(\mathbb{C})$,
define the  {\it McKay graph}
$\tilde{\Delta}_\Gamma$ as follows:
the vertices are $0,1,2,\ldots,\ell$ which
index the
irreducible representations
$R_0,R_1,\ldots,R_\ell$ of $\Gamma$.

Let $Q$ be the 2-dimensional $\Gamma$ coming
from
$\Gamma<\text{SL}_2(\mathbb{C})$.
(Note that $Q$ is one of the $R_i$'s unless
$\Gamma=\mathcal{C}_m$,
in which case $Q$ is reducible.)
The number of edges between $i$ and $j$ in
$\tilde{\Delta}_\Gamma$
is the multiplicity of $R_j$ in $Q \otimes
R_i$.
So: take the representation $R_i$, tensor it
with the
2-dimensional representation $Q$,
express it as a direct sum of irreducible
representations
and perhaps $R_j$ will appear with nonzero
multiplicity,
in which case you put $R_j$ with that number
of edges.
So in general it's a multigraph, though in
most
cases its a simple graph.
Also note that this multiplicity coincides
with the multiplicity of $R_i$ in $Q \otimes
R_j$
since $Q^* \cong Q$; that is the graph is not
directed.

The surprising thing is that when we do this
we obtain
the affine Dynkin diagrams. McKay's famous
observation says:
\begin{enumerate}
\item $\tilde{\Delta}_\Gamma$ is a
simply-laced affine Dynkin
diagram.

\item This is a bijection:
$$
\left\{\substack{\text{nontrvial finite} \\
\Gamma \subset
\text{SL}_2(\mathbb{C})}\right\}\Big/
\text{conj.}
\leftrightarrow
\left\{\substack{\text{simply-laced} \\
\text{affine Dynkin diagrams}}\right\}.
$$
\end{enumerate}

\noindent
The correspondance reads
\begin{align*}
\mathcal{C}_m &\xymatrix{\ar@{<->}[r]&}
\widetilde{A}_{m-1}\\
\mathcal{D}_m &\xymatrix{\ar@{<->}[r]&}
\widetilde{D}_{m+2}\\
\mathcal{T}&\xymatrix{\ar@{<->}[r]&}\tilde{E}
_6\\
\mathcal{O}&\xymatrix{\ar@{<->}[r]&}\tilde{E}
_7\\
\mathcal{I}&\xymatrix{\ar@{<->}[r]&}\tilde{E}
_8.
\end{align*}

\medskip\noindent
Now we explain the geometric approach.
Given a finite $\Gamma \subset
\text{SL}_2(\mathbb{C})$, consider
the quotient $Y_\Gamma=\Gamma\setminus Q$ as
an algebraic variety,
that is, (since $Q$ is just a representation
of $\Gamma$, i.e.
a 2-dimensional affine space with an action
of $\Gamma$)
we can take its quotient in the sense
described before:
$Y_\Gamma$ is the affine variety with
coordinate algebra
$\mathbb{C}[x,y]^\Gamma$.

\begin{lemma}
\label{lemma-action-is-free-outside-origin}
The action of $\Gamma$ is free on
$Q\setminus\{0\}$ is free.
\end{lemma}



\section{Borel-Weil-Bott theorem}
\label{section-Borel-Weil-Bott-thoerem}

[These are {\it very} sketchy notes after a
conversation with Lyalya Guseva on
Borel-Weil-Bott theorem]

Borel-Weil-Bott theorem is a device for
computing cohomologies of sheaves.

Take the Grassmanian $\text{Gr}(n,k)$, which
may be obtained as a quotient $G$
of a group $G$ by a parabolic subgroup $P$.
We would like it if $P$ was a
semisimple group since those are classified,
but unfortunately it is not.
So instead we use so-called Levi quotient
denoted by $L$ which is semisimple
and allows us to understand the cohomologies
of the variety $G/P$.

Fortunately the representation theory of
$\text{Gr}(k,n)$ is well known, in fact
we get $\text{GL}_k \times \text{GL}_{n-k}$.

There are functors associated to $L$, which
are described by a sequence of
numbers $a_1,\ldots,a_n$. Along with other
sequences of numbers
$k_1,\ldots,k_\ell$, and $\rho$ (the latter
is a concept in representation
theory but we may ultimately think of it as
another sequence of numbers) we may
construct an action of the symmetric group
$S^n$ acting on these sequences and
obtain a result concerning the cohomology
$H^{\ell(\sigma)}(L)$, and in
particular we find that if two entries in our
list of numbers coincide, the
cohomology will vanish.


\bibliography{my}
\bibliographystyle{amsalpha}

\end{document}
